\exnumber{03}
\extitle{UX II - Ex 4 - Énoncé du problème}
\exfiles{ex03/PDF, nomme Name\_FirstName\_DeliverableName\_Date\_VersionNumber}
\exauthorize{Ceux nommés dans l'exercice}
\makeheaderfiles

\textbf{Instructions :}

À partir de votre travail précédent et de votre empathy map, définissez un problem statement qui clarifierait le problème principal que vous voulez traiter.

Vous pouvez utiliser l'exercice How Might We en binôme.

\textbf{Livrable :}

\begin{itemize}
\item Un document PDF avec un énoncé de problème commençant par un « How Might We »
\item Déposez le PDF sur Github, nommé ainsi : « Name\_FirstName\_DeliverableName\_Date\_VersionNumber ». Exemple : « Doe\_John\_MidFiPrototype\_05242024\_V1 ».
\item Ajoutez un titre dans le document.
\item Soignez l'aspect global du document. Il doit être clair, simple et lisible. Vous pouvez regarder canva.com, sans surcharger le dessin de détails.
\end{itemize}

\textbf{Ressources :}

\begin{itemize}
\item Bloqué sur un problème ? Demandez simplement How Might We \url{https://relab.academy/design-thinking/stuck-on-a-problem-just-ask-how-might-we/}
\item Problem statement \url{https://www.indeed.com/career-advice/career-development/what-is-a-problem-statement}
\item Utiliser les questions « How Might We » pour idéer sur les bons problèmes \url{https://www.nngroup.com/articles/how-might-we-questions/}
\end{itemize}

\textbf{Citation :}

« Du désordre, tirez la simplicité. De la discorde, trouvez l'harmonie. Au milieu de la difficulté se trouve l'opportunité. » Albert Einstein

